% Extracto íntegro por residencia del cierre físico-matemático sucesor.
\section{Réponse auto-inertielle et principe de Mach à partir de l'Ambivalence}
\label{sec:hmtf2-mach-ambivalencia}

Soit \(b:\mathcal X_\infty^{\rm enr}\to B_\infty\) la trace globale de
frontière. Sur \(B_\infty\) sont conservés les lecteurs positifs
\(A(\beta)\), \(V(\beta)\), l'horloge globale \(Q(\beta)>0\) et la mémoire
\(T(\beta)\). La charnière d'Ambivalence fournit
\[
 r_{\rm av}=\frac{\pi_{\rm HMT}}{\varphi_{\rm HMT}^2},
 \qquad
 r_{\rm av}^J=\frac{\varphi_{\rm HMT}}{\pi_{\rm HMT}^2}.
\]
Nous définissons des poids, sans données gravitationnelles extérieures,
\[
 w_A(\beta)
 =\frac{A(\beta)r_{\rm av}}
 {A(\beta)r_{\rm av}+V(\beta)r_{\rm av}^J},
 \qquad
 w_V(\beta)=1-w_A(\beta).
\]
Soient \(P_A,P_V\) des projections orthogonales complémentaires de l'espace local
de réponse. Pour des énergies internes \(E_1,E_2\), la réponse
auto-inertielle HMT est
\[
 \overline I_v^{\rm HMT}(\beta)
 =\frac{E_1E_2}{Q(\beta)^2}
 \bigl(w_A(\beta)P_A+w_V(\beta)P_V\bigr),
\]
et
\[
 I_v^{\rm HMT}
 =\overline I_v^{\rm HMT}\circ b.
\]

\begin{theorem}[Clôture machienne de la réponse d'Ambivalence]
\label{thm:hmtf2-respuesta-mach}
La réponse \(I_v^{\rm HMT}\) est autoadjointe, positive et constante sur chaque
fibre de \(b\). Elle se factorise donc de manière unique par la totalité globale.
L'involution d'Ambivalence qui échange
\[
 (A,r_{\rm av},P_A)\longleftrightarrow
 (V,r_{\rm av}^J,P_V)
\]
conjugue \(\overline I_v^{\rm HMT}\) sans changer son spectre. S'il existe
\(x,x'\) ayant la même restriction locale, \(b(x)\ne b(x')\) et un
\(Q,A\) ou un \(V\) différent, alors la réponse ne se factorise pas uniquement par la restriction
locale et elle est machienne au sens fort.
\end{theorem}

\begin{proof}
Les poids sont positifs et leur somme vaut un ; la combinaison de projections est
positive et autoadjointe. La formule ne dépend que de \(b(x)\), elle est donc
constante sur ses fibres et le critère de factorisation par la
totalité s'applique. L'échange permute les deux blocs spectraux. La dernière
affirmation est précisément le contrôle négatif face à une réponse
purement locale.
\end{proof}


La construction possède en outre une lecture géométrique retrouvée. Pour l'
état conjoint \(\Gamma=(x,\mathcal O,m)\), soit
\(\nabla^{\rm tot}_{b(x)}\) la connexion induite par la frontière et soit
\(\pi_S\) la projection sur le sous-système. Si
\[
 \nabla^{\rm tot}_{\dot\Gamma}\dot\Gamma=0,
\]
l'accélération publiée est
\[
 a_S
 =\nabla^S_{\dot\gamma_S}\dot\gamma_S
 =\mathcal K_{\pi_S}^{\,b(x)}(\dot\Gamma,\dot\Gamma).
\]
Comme le défaut dépend de la connexion induite par \(b(x)\), il se
factorise aussi par la totalité. Cela clôt mathématiquement le remplacement de la
dichotomie inertiel/non inertiel par des systèmes auto-inertiels : l'état total
est géodésique et l'accélération locale est un défaut de projection. L'égalité de
cette réponse à une accélération gravitationnelle observée relève de la
\textsc{reconnaissance extérieure}.

